\section*{Anexo I. Glosario de términos}

A continuación, se definen los conceptos técnicos, acrónimos y vocabulario especializado utilizados a lo largo de este Trabajo de Fin de Grado, con el objetivo de facilitar la lectura y comprensión del documento.

\textbf{CAPEX (Capital Expenditure):} 
Gastos de capital en bienes físicos, inversión inicial en infraestructura.

\textbf{CSRF (Cross-Site Request Forgery):} 
Vulnerabilidad de seguridad web que permite a un atacante inducir a los usuarios a realizar acciones que no pretenden realizar en una aplicación en la que están autenticados.

\textbf{Data Leakage (Fuga de datos):} 
Error en el aprendizaje automático que ocurre cuando se utiliza información externa al conjunto de entrenamiento para crear el modelo, lo que resulta en métricas de rendimiento irrealmente optimistas que no se generalizan a datos nuevos.

\textbf{ETL (Extract, Transform, Load):} 
Proceso de integración de datos que consiste en extraer datos de diversas fuentes, transformarlos para cumplir con requisitos de negocio o técnicos, y cargarlos en un sistema de destino.

\textbf{EWMA (Exponentially Weighted Moving Average):} 
Método estadístico para series temporales que aplica pesos decrecientes a los datos históricos, priorizando la información más reciente.

\textbf{Feature Engineering (Ingeniería de características):} 
Proceso de transformar datos en bruto en variables que representen mejor el problema subyacente a los algoritmos predictivos, mejorando así la precisión del modelo.

\textbf{GridSearch:} 
Técnica de optimización de hiperparámetros que realiza una búsqueda exhaustiva a través de un subconjunto especificado manualmente del espacio de parámetros de un algoritmo de aprendizaje automático.

\textbf{H2H (Head-to-Head):} 
Comparación directa entre dos entidades o modelos bajo las mismas condiciones.

\textbf{JWT (JSON Web Token):} 
Estándar abierto basado en JSON para la creación de tokens de acceso que permiten la propagación de identidad y privilegios de forma segura entre un cliente y un servidor.

\textbf{Kanban:} 
Metodología ágil para la gestión del flujo de trabajo que se centra en la visualización del trabajo, la limitación del trabajo en curso (WIP) y la mejora continua del proceso de entrega.

\textbf{MVC (Model-View-Controller):} 
Patrón de arquitectura de software que separa la lógica de negocio (Modelo), la interfaz de usuario (Vista) y la gestión de eventos (Controlador) para facilitar el mantenimiento y la escalabilidad.

\textbf{OPEX (Operational Expenditure):} 
Gastos operativos recurrentes, asociado comúnmente al modelo de pago por uso de servicios, mantenimiento o salarios.

\textbf{Outliers (Valores atípicos):} 
Observaciones en un conjunto de datos que se encuentran a una distancia anormal de otros valores y que pueden indicar variabilidad en la medición, errores experimentales o novedades de interés.

\textbf{Prop Drilling:} 
Fenómeno en frameworks de desarrollo frontend donde los datos se pasan a través de múltiples niveles de componentes que no los necesitan, solo para llegar a un componente anidado.

\textbf{Proxy:} 
Servidor o programa que actúa como intermediario entre las peticiones de recursos de un cliente y el servidor de destino, utilizado a menudo para seguridad, filtrado o mejora del rendimiento.

\textbf{Random Forest:} 
Algoritmo de aprendizaje supervisado compuesto por una multitud de árboles de decisión, cuya salida es la clase que representa la moda de las clases o la media de las predicciones de los árboles individuales.

\textbf{Scraping:} 
Técnica automatizada para extraer grandes cantidades de información de sitios web de forma estructurada para su posterior análisis o almacenamiento.

\textbf{SHAP (SHapley Additive exPlanations):} 
Método basado en la teoría de juegos para explicar el resultado de modelos de aprendizaje automático, asignando a cada característica un valor de importancia para una predicción específica.

\textbf{XAI (Explainable Artificial Intelligence):} 
Conjunto de procesos y métodos que permiten a los usuarios humanos comprender y confiar en los resultados y predicciones creados por algoritmos de aprendizaje automático.

\textbf{XGBoost (eXtreme Gradient Boosting):} 
Implementación optimizada de algoritmos de aumento de gradiente diseñada para ser altamente eficiente, flexible y portátil, siendo uno de los algoritmos más utilizados en ciencia de datos.